% ============================================================================
% APPD.TEX - Lampiran D: Best Practices \& Tips
% ============================================================================

\chapter{Best Practices \& Tips}

\section{Code Best Practices}

\subsection{Performance}

\begin{itemize}
\item Cache GetComponent() di Start/Awake
\item Hindari operations mahal di Update()
\item Gunakan object pooling untuk Instantiate/Destroy
\item Minimize string operations
\item Use appropriate data structures
\end{itemize}

\subsection{Organization}

\begin{itemize}
\item Organize scripts dalam folder yang jelas
\item Gunakan namespaces untuk grouping
\item Comment code yang kompleks
\item Follow naming conventions
\item Keep functions focused dan kecil
\end{itemize}

\subsection{Debugging}

\begin{itemize}
\item Gunakan Debug.Log dengan bijak
\item Remove debug logs sebelum release
\item Use breakpoints untuk debugging kompleks
\item Check Console untuk warnings
\end{itemize}

\section{Unity Best Practices}

\subsection{Project Organization}

\begin{itemize}
\item Organize assets dalam folder (Scripts, Art, Audio, Prefabs)
\item Use consistent naming conventions
\item Keep scene hierarchy organized
\item Use empty GameObjects sebagai organizers
\end{itemize}

\subsection{Performance}

\begin{itemize}
\item Use LOD untuk models kompleks
\item Optimize textures (compression, size)
\item Use object pooling
\item Minimize draw calls (batching)
\item Use occlusion culling
\end{itemize}

\subsection{Workflow}

\begin{itemize}
\item Save scene secara teratur
\item Use version control (Git)
\item Test di target platform
\item Document important decisions
\item Keep backups
\end{itemize}

\section{Common Mistakes}

\begin{itemize}
\item \textbf{GetComponent di Update}: Cache di Start
\item \textbf{Instantiate/Destroy setiap frame}: Use object pooling
\item \textbf{String concatenation di Update}: Update hanya saat berubah
\item \textbf{Tidak unsubscribe events}: Unsubscribe di OnDisable
\item \textbf{Forgetting to set Time.timeScale back}: Reset setelah pause
\end{itemize}

\section{Tips untuk Pemula}

\begin{enumerate}
\item Start dengan project kecil
\item Learn dari tutorial dan dokumentasi
\item Experiment dengan fitur Unity
\item Join komunitas Unity
\item Practice regularly
\item Don't be afraid to make mistakes
\item Read error messages carefully
\item Use Profiler untuk optimasi
\end{enumerate}

\section{Resources untuk Belajar}

\begin{itemize}
\item Unity Learn: \url{https://learn.unity.com/}
\item Unity Documentation: \url{https://docs.unity3d.com/}
\item Unity Forum: \url{https://forum.unity.com/}
\item Unity Answers: \url{https://answers.unity.com/}
\item YouTube Channels: Brackeys, Code Monkey, dll
\end{itemize}

\section{Kesimpulan}

Pengembangan game adalah proses pembelajaran yang berkelanjutan. Teruslah belajar, experiment, dan berbagi dengan komunitas. Practice makes perfect!
